\chapter{Exact results and limiting-case validation}
\label{chap:validation}

A reference solver earns trust by failing loudly when any layer disagrees with
an independent result.  The validation suite is deliberately redundant:
operator algebra, matrix properties, analytic spectra, limiting cases, and
iterative residuals constrain different classes of mistakes.

\section{The validation ladder}

\begin{center}
\begin{tabular}{p{0.24\textwidth}p{0.66\textwidth}}
\toprule
Check & What it detects \\
\midrule
Basis dimension & Missing, duplicated, or out-of-sector states \\
Anticommutators & Global-ordering and parity errors \\
Hermiticity & Missing reverse hops and inconsistent signs \\
Two-site formula & Interaction normalization and hopping scale \\
Full $U=0$ spectrum & Bond convention, degeneracies, and spin factorization \\
Atomic limit & Double-occupancy diagonal \\
Large-$U$ gap & Low-energy scale and excited-state ordering \\
CSR versus matrix-free & Divergent implementations \\
Residual norms & Eigensolver convergence and returned-vector integrity \\
\bottomrule
\end{tabular}
\end{center}

\section{Half-filled two-site model}

For $L=2$, $N_\uparrow=N_\downarrow=1$, and one unique bond, the ground-state
energy is
\begin{equation}
E_0=\frac{U}{2}-\frac{1}{2}\sqrt{U^2+16t^2}.
\label{eq:twosite}
\end{equation}
This single formula tests the relative factors in both Hamiltonian terms.  It
should be checked at several positive and negative $U$ values and more than one
$t$.

At large repulsive $U$, the triplet lies at zero energy while the singlet is
lowered, giving
\begin{equation}
\Delta_{ST}=-E_0
=\frac{\sqrt{U^2+16t^2}-U}{2}
=\frac{4t^2}{U}+O(U^{-3}).
\end{equation}
Thus the solver recovers the superexchange scale $J=4t^2/U$.

\section{Noninteracting limit}

At $U=0$, diagonalize the one-particle hopping problem independently.  For an
open chain,
\begin{equation}
\epsilon_m=-2t\cos\left(\frac{m\pi}{L+1}\right),
\qquad m=1,\ldots,L.
\end{equation}
For a periodic chain with $L\ge3$,
\begin{equation}
\epsilon_m=-2t\cos\left(\frac{2\pi m}{L}\right),
\qquad m=0,\ldots,L-1.
\end{equation}
Under the unique-edge convention, $L=2$ instead has energies $(-|t|,|t|)$.

Every many-body energy in the fixed sector is a sum of occupied one-particle
energies,
\begin{equation}
E=\sum_{a\in A_\uparrow}\epsilon_a
 +\sum_{b\in A_\downarrow}\epsilon_b.
\end{equation}
Comparing the complete small-system spectrum---including repeated values---is
stronger than checking only the ground energy.  Degeneracies are part of the
expected answer, not noise to be discarded.

\section{Atomic limit and implementation agreement}

At $t=0$, every basis state is already an eigenstate and
\begin{equation}
E(s)=U\,\operatorname{popcount}(b_\uparrow\mathbin{\texttt{\&}}b_\downarrow).
\end{equation}
There must be no off-diagonal matrix elements.  Finally, for seeded random
vectors, require
\begin{equation}
\norm{H_{\mathrm{CSR}}x-H_{\mathrm{free}}x}_2
\end{equation}
to be consistent with floating-point roundoff, and require Eq.~\eqref{eq:residual}
for every computed eigenpair.

Run the complete executable validation with
\begin{verbatim}
.venv/bin/python -m pytest -q
\end{verbatim}
The tutorial programs supplement rather than replace these automated tests.

\paragraph{Checkpoint.}
A new representation should reproduce this entire ladder on the largest
shared system size before its performance or scaling claims are considered.
